\begin{table}[htbp]
\centering
\small
\caption{Шесть инвариантов допуска к расчёту серии v2}
\label{tab:s11}
\begin{tabular}{>{\RaggedRight\hspace{0pt}\arraybackslash}p{0.302\linewidth}>{\RaggedRight\hspace{0pt}\arraybackslash}p{0.302\linewidth}>{\RaggedRight\hspace{0pt}\arraybackslash}p{0.302\linewidth}}
\toprule
№ & Инвариант & Результат \\
\midrule
1 & реестр содержит 1882 документа & пройден \\
2 & ни один из 34 исключённых не попал в матрицу & пройден \\
3 & raw и normalized 1814 неизменённых совпадают с v4 & пройден, 114 282 сверенных значения, расхождений ноль \\
4 & percentile-реализация воспроизводит v4 на составе v4 & пройден, 88 114 значений \\
5 & перцентили v5 посчитаны на новой референсной выборке & пройден \\
6 & результаты серии v1 не перезаписаны & пройден \\
\bottomrule
\end{tabular}
\begin{minipage}{\linewidth}\footnotesize\vspace{2pt}
\textit{Примечание.} Единица анализа — инвариант. Знаменатель: шесть инвариантов; допуск выставляется только при шести пройденных. Данные: манифест допуска, ревизия 6. Таблица описывает шлюз; оценок в ней нет. Шкала: шкалы нет. Сокращения: raw и normalized — сырое и нормированное значение признака; percentile-реализация — расчёт перцентилей на референсной выборке. Пропуски: пропусков нет. Поправка на множественные сравнения не применялась.
\end{minipage}
\end{table}
